
\documentclass[11pt,a4paper]{moderncv}
\moderncvtheme[green]{classic}                
\usepackage[T1]{fontenc}
\usepackage[utf8x]{inputenc}
\usepackage[italian]{babel}
\usepackage[scale=0.8]{geometry}
\recomputelengths                             

\fancyfoot{} 

\renewcommand*{\namefont}{\fontsize{25}{28}\mdseries\upshape}

\firstname{Patricio Nicolas}
\familyname{Massaro Rocca}

\mobile{+1-202-754-6052}                    
\email{patomassaro@gmail.com}                

\begin{document}
\maketitle
\vspace{-10mm}

\section{Work experience}

\cventry{2020 - 2024 }{Lead Machine Learning Engineer}{\href{https://www.bairesdev.com/}{Bairesdev}}{}{}{
\begin{itemize}
    \item Pioneered the Machine Learning department, leading a team to train and launch machine learning models for client segmentation and rating, using AWS, Databricks, Airflow, MLFlow, and SparkML, which led to generating 500 million daily predictions and a 100\% increase in marketing campaign conversion rates.
    \item Developed and integrated machine learning algorithms within production systems to reduce spam rates by 30\%, enhancing overall platform reliability and security.
    \item Implemented a natural language processing deep learning model to classify inbound messages, reducing human workload by 37\%.
    \item Optimized data architecture to enhance ML models training, increasing performance metrics by 10\%.
    \item Utilized Spark SQL and Tableau for data analysis to inform strategic decisions among top management.
\end{itemize} 
}

\cventry{2023-2024}{Thesis working student}{\href{https://ororatech.com/}{Ororatech}}{}{}{  
\begin{itemize}
    \item Trained GANs in Pytorch for remote-sensing blind super-resolution, achieving a 20\% error reduction in land surface temperature estimation with applications in wildfire management.
    \item Developed a deep convolutional model for fine-scale georeferencing of thermal data products, improving accuracy in disaster monitoring applications.
    \item Published paper at IGARSS 2023 on multi-spectral remote sensing super-resolution.
\end{itemize} }

\cventry{ 2017 - 2020}{Business Analytics Engineer}{\href{https://www.tenaris.com/en}{Tenaris}}{}{}{
 \begin{itemize}
     \item Led the architectural design and workflow development for migrating the entire organization to a centralized cloud-based PowerBI application for 30,000 contributors.
     \item Managed the population of the data warehouse using Synapse and Data Factory and designed the semantic layer using Azure Analysis Services, reducing costs by 17\%.
     \item Defined platform adoption KPIs and action plans for their improvement.
 \end{itemize}
 }

\cventry{ 2017 - 2018}{Research Intern}{University of Buenos Aires}{}{}{ 
\begin{itemize}
    \item Authored two publications on the effects of piezoelectric devices in optoacoustic image reconstruction.
\end{itemize}}

\cventry{2014 - 2017}{Systems Administrator}{Masstec S.R.L}{}{}{
    \begin{itemize}
        \item Administered server, workstation, and networking systems for small businesses.
    \end{itemize}
}

\cventry{2013- 2014}{Network Administrator}{Education Ministry, Buenos Aires}{}{}{
\begin{itemize}
    \item Coordinated the distribution of laptops and integrated ICT into educational practices under the "Conectar Igualdad" program, enhancing digital literacy for over 250 students. 
\end{itemize}
}

\section{Education}
\cventry{2021-2023}{Elite MSc Data Science}{\href{https://www.m-datascience.mathematik-informatik-statistik.uni-muenchen.de/index.html}{Ludwig-Maximilians-Universität}}{Munich}{Germany}{Thesis: Blind Super Resolution Applied to Thermal Remote Sensing Data Products}
\cventry{2020-2021}{MSc Data Mining }{\href{http://datamining.dc.uba.ar/datamining/}{Universidad de Buenos Aires}}{}{Argentina}{}
\cventry{2010-2018}{Electronic Engineering}{\href{https://www.fi.uba.ar/}{Universidad de Buenos Aires}}{}{Argentina}{Thesis: Improvements in optoacoustic image reconstruction using deconvolution filters.} 

\section{Language}
\cvline{Spanish}{Native language.}
\cvline{English}{Bilingual proficiency.}
\cvline{German}{Professional working proficiency.}

\vspace{-3mm}

\end{document}
